\section{Pointwise allocation of model error}\label{sec:allocation}

The construction uses an upper bound on local model error. A necessary
partition-size bound requires additional information: the model must lose
a specified amount in the interior of a box. This appendix gives such a
bound. It also explains why a single tolerance share for each bag can be
too restrictive when different near-optimal points activate different
bags. No quadratic-growth assumption is used here.

\subsection{A lower gap contract and variable-radius coverings}
For a bag box $B=\prod_{i\in V_t}[\ell_i^B,u_i^B]$, write
\[
 q_{B,K}(z)=\sum_{i\in K}(z_i-\ell_i^B)(u_i^B-z_i)
 \qquad(K\subseteq V_t,\ z\in B).
\]
Fix $\alpha>0$ and coordinate sets $K_t\subseteq V_t$. The additional
\emph{lower gap contract} is
\begin{equation}\label{eq:allocation-gap}
 a_t(z)-m_{t,B}(z)\ge\alpha q_{B,K_t}(z)
 \qquad(z\in B).
\end{equation}
It is imposed on every model used by a certificate under consideration.
It does not follow from \eqref{eq:model-contract}: an exact convex model
has zero error even in the interior. Thus the results below do not apply
to the main construction solely on the basis of its upper error bound.

Let
\[
 E(\eta)=\{x\in X:F(x)-f^*\le\eta\},\qquad \eta\ge0.
\]
For a nonnegative function $e:X_{V_t}\to[0,\infty)$ and a set
$A\subseteq X_{V_t}$, let $\mathcal N_{t,2}(e;A)$ be the least cardinality
of a finite set $C\subseteq\R^{K_t}$ such that
\[
 \text{for every }z\in A\text{ there is }v\in C\text{ with }
 \norm{z_{K_t}-v}\le\sqrt{e(z)/\alpha}.
\]
If no finite set exists, the number is $+\infty$.
Define $\mathcal N_{t,\infty}(e;A)$ by replacing the Euclidean norm with
the sup norm. Empty coordinate sets have the usual one-point space.
Since a Euclidean ball is contained in the sup-norm ball of the same
radius, $\mathcal N_{t,2}\ge\mathcal N_{t,\infty}$.

A family of nonnegative functions $(e_t)$ is \emph{admissible at
$\eta$} if
\begin{equation}\label{eq:allocation-admissible}
 \sum_t e_t(x_{V_t})\le\eps+\eta\qquad(x\in E(\eta)).
\end{equation}
For $\nu\in\{2,\infty\}$, set
\begin{equation}\label{eq:allocation-phi}
 \Phi_\nu(\eps,\eta)=
 \inf_{(e_t)\text{ admissible at }\eta}
 \sum_t2^{-|K_t|}\mathcal N_{t,\nu}
       (e_t;\operatorname{proj}_{V_t}E(\eta)).
\end{equation}
The functions can depend on the point in the bag. The infimum imposes no
regularity or computability requirement; it defines a necessary size
bound, rather than an algorithm for locating a partition.

\begin{theorem}[Pointwise Euclidean covering bound]
\label{thm:allocation-pointwise}
Let a separator certificate satisfy \eqref{eq:allocation-gap} and
$f^*-\LB\le\eps$. For every bag define
\[
 Q_t(z)=\min_{B\in\mathcal L_t:\ z\in B}q_{B,K_t}(z).
\]
Then $(\alpha Q_t)$ is admissible at every $\eta\ge0$, and
\begin{equation}\label{eq:allocation-leaf-bound}
 |\mathcal L_t|\ge2^{-|K_t|}
 \mathcal N_{t,2}(\alpha Q_t;\operatorname{proj}_{V_t}E(\eta))
 \ge2^{-|K_t|}
 \mathcal N_{t,\infty}(\alpha Q_t;\operatorname{proj}_{V_t}E(\eta)).
\end{equation}
Consequently the certificate size satisfies
\begin{equation}\label{eq:allocation-size-bound}
 \mathcal S\ge\sup_{\eta\ge0}\Phi_2(\eps,\eta)
 \ge\sup_{\eta\ge0}\Phi_\infty(\eps,\eta).
\end{equation}
\end{theorem}
\begin{proof}
The minimum defining $Q_t$ exists because the leaf family is a finite
cover. Fix $x\in E(\eta)$ and, at each bag, choose a containing leaf
$B_t$ attaining $Q_t(x_{V_t})$. Cancellation along the consistent point,
\eqref{eq:model-chain}, and the lower gap contract give
\[
 \alpha\sum_tQ_t(x_{V_t})
 \le\sum_t\bigl[a_t(x_{V_t})-m_{t,B_t}(x_{V_t})\bigr]
 \le F(x)-\LB\le\eta+\eps.
\]
This proves admissibility, including at points on partition boundaries.

For $z\in X_{V_t}$ choose a leaf $B$ attaining $Q_t(z)$. If $d_i$ is
the distance from $z_i$ to the nearer endpoint of $[\ell_i^B,u_i^B]$,
then
$d_i^2\le(z_i-\ell_i^B)(u_i^B-z_i)$: the smaller of two nonnegative
numbers has square at most their product. Choosing the nearer endpoint
in every coordinate of $K_t$ gives a vertex $v$ of $B_{K_t}$ with
\[
 \norm{z_{K_t}-v}^2=\sum_{i\in K_t}d_i^2\le Q_t(z).
\]
The union of the projected vertices of all leaves has at most
$2^{|K_t|}|\mathcal L_t|$ points and is therefore a permissible covering
set. This proves \eqref{eq:allocation-leaf-bound}. Sum over bags, use
$\mathcal S\ge\sum_t|\mathcal L_t|$, and take the infimum in
\eqref{eq:allocation-phi} and then the supremum over $\eta$.
\end{proof}

For each fixed $\eta$, admissibility also implies
$e_t(z)\le\eps+\eta$ on $\operatorname{proj}_{V_t}E(\eta)$: lift $z$ to
a point of $E(\eta)$ and use nonnegativity of the other terms.
Thus \eqref{eq:allocation-size-bound} contains the corresponding
constant-radius projected covering bound. It preserves the stronger
pointwise restriction on the sum of the radii squared.

\subsection{Different points can use different bags}
\label{sec:allocation-staircase}
The following finite-state example isolates the allocation issue. Its
separator variables are binary; it is not an instance of the continuous
box model used in the regridding theorem. The certificate inequalities
retain the same meaning on each continuous slice, with singleton cells
for binary separator values. Affine functions on these singleton cells
can be constants.

For a bounded set $A$ in a coordinate space, let $\mathcal C(A,r)$ be
the least number of closed axis-parallel cubes of side length $r$ covering
$A$. Unlike the radius convention above, $r$ here is the full side
length. A measure based on constant positive shares is
\begin{equation}\label{eq:allocation-constant-shares}
 \Psi(\eps)=\inf_{\substack{\eps_t>0\\\sum_t\eps_t=\eps}}
 \sum_t\sup_{\eta\ge0}\mathcal C\left(
 \operatorname{proj}_{V_t}E(\eta),
 \sqrt{(\eps_t+\eta)/\alpha}\right).
\end{equation}

\begin{proposition}[A staircase of flat slices]
\label{prop:allocation-staircase}
Fix $K\ge2$, $d\ge1$, $0<\alpha\le1$, $0<\eps<1$, and
$P\ge1+\alpha d/4$. Let $s_1,\ldots,s_{K+1}\in\{0,1\}$ and
$y_t\in[0,1]^d$, and use the path bags
$V_t=\{s_t,y_t,s_{t+1}\}$ with costs
\[
 a_t=(1-s_{t+1}+s_t)\norm{y_t}^2+
       P\max\{s_t-s_{t+1},0\}^2.
\]
On a box $B$ in a fixed binary slice supply the convex model
$m_{t,B}=a_t-\alpha q_{B_y,\{1,\ldots,d\}}$.
Then $f^*=0$ and
\begin{equation}\label{eq:allocation-staircase-lower}
 \Psi(\eps)\ge K(\alpha K/\eps)^{d/2}.
\end{equation}
There is a finite-state separator certificate, using zero slopes, with
$\UB=0$, $\LB\ge-\eps$, and at most
\begin{equation}\label{eq:allocation-staircase-upper}
 K\left[\ceil{\sqrt{\alpha d/(2\eps)}}^d+
 2(J+1)(4/\theta)^d+1\right]+2(K-1)
\end{equation}
leaves and cells, where
\[
 h_0=\sqrt{\frac{2\eps}{K d\alpha^2}},\quad
 J=\max\{0,\ceil{\log_2(1/h_0)}\},\quad
 \theta=\max\{2^{-j}:j\in\N\cup\{0\},\
                 2^{-j}\le\min(1,2/\sqrt{\alpha d})\}.
\]
For fixed $d,\alpha,\eps$, the ratio of \eqref{eq:allocation-staircase-lower}
to the size in \eqref{eq:allocation-staircase-upper} is at least
$cK^{d/2}/(1+\log K)$ for a constant $c>0$ independent of $K$.
\end{proposition}
\begin{proof}
Every cost is nonnegative on the binary domain. A zero-cost point has no
descent in $s$, has $y_t=0$ at every equal pair, and can have arbitrary
$y_t$ at its ascent, if any. In particular, for each $t$, $E(0)$ contains
the slice with $s_1=\cdots=s_t=0$, $s_{t+1}=\cdots=s_{K+1}=1$,
$y_t\in[0,1]^d$, and all other $y_j=0$. A cube of side $r$ meets the
projection of this slice in $d$-dimensional volume at most $r^d$.
Consequently the $t$th term of \eqref{eq:allocation-constant-shares} is
at least $(\alpha/\eps_t)^{d/2}$, by taking $\eta=0$. Convexity of
$u\mapsto u^{-d/2}$ and $\sum_t\eps_t=\eps$ prove
\eqref{eq:allocation-staircase-lower}.

Partition each bag separately on its four binary slices. On the ascent
slice $(0,1)$ use a uniform grid with
$n_1=\ceil{\sqrt{\alpha d/(2\eps)}}$ intervals per continuous coordinate.
Its $n_1^d$ leaves have model minimum at least
$-\alpha d/(4n_1^2)\ge-\eps/2$.

On each equal slice use the following partition. Its central cube is
$[0,2^{-J}]^d$. For $j=0,\ldots,J-1$, partition
$[0,2^{-j}]^d\setminus[0,2^{-(j+1)})^d$ by retaining the cubes outside
the inner cube from the uniform grid of side
$\theta2^{-(j+1)}$ on the outer cube. The dyadic choice of $\theta$
makes both boundaries grid boundaries. These closed cubes and the central
cube cover $[0,1]^d$ and have disjoint interiors. There are at most
$1+J(2/\theta)^d\le(J+1)(4/\theta)^d$ cubes. Each noncentral cube $B$
satisfies
$\width(B)\le\theta\dist_\infty(B,0)$, since at least one of its lower
coordinates is at least $2^{-(j+1)}$. Its model is therefore bounded below
by
\[
 \dist_\infty(B,0)^2-\frac{\alpha d}{4}\width(B)^2
 \ge\dist_\infty(B,0)^2(1-\alpha d\theta^2/4)\ge0.
\]
On the central cube of side $h=2^{-J}\le\min(h_0,1)$, the minimum is
\[
 \min_{[0,h]^d}\left[(1+\alpha)\norm y^2-
                          \alpha h\sum_i y_i\right]
 =-\frac{d\alpha^2h^2}{4(1+\alpha)}\ge-\frac\eps{2K}.
\]
On the descent slice $(1,0)$ use one leaf. Its model minimum is at least
$P-\alpha d/4\ge1$. The models are convex on all four slices: the
coefficient of $\norm y^2$ is $1-\sigma'+\sigma+\alpha>0$.

To construct the affine certificate, let $\mu_t(\sigma,\sigma')$ be
the least model minimum among the leaves in slice $(\sigma,\sigma')$.
Put $b_{K+1}(\sigma)=0$ and recursively set
\[
 b_t(\sigma)=\min_{\sigma'\in\{0,1\}}
                  [\mu_t(\sigma,\sigma')+b_{t+1}(\sigma')].
\]
For each nonroot separator use its two singleton cells, with constant
minorants $b_t(0),b_t(1)$. On a parent leaf the child bound is the constant
$b_{t+1}(\sigma')$ determined by its fixed right state. The child and
local inequalities hold by this recurrence; set
$\LB=\min_{\sigma\in\{0,1\}}b_1(\sigma)$ at the root.
Equivalently, $\LB$ is the minimum of $\sum_t\mu_t(s_t,s_{t+1})$ over
binary sequences. If a sequence has $A_s$ ascents and $D_s$ descents,
then $A_s-D_s=s_{K+1}-s_1\le1$. Its value is at least
\[
 -A_s\eps/2-K\eps/(2K)+D_s
 \ge-\eps+D_s(1-\eps/2)\ge-\eps.
\]
The all-zero point gives $\UB=0$. Counting the four slice partitions and
the two cells at each separator yields
\eqref{eq:allocation-staircase-upper}. For fixed $d,\alpha,\eps$, its
size is $O(K(1+\log K))$, while
\eqref{eq:allocation-staircase-lower} grows as $K^{1+d/2}$, proving the
last assertion.
\end{proof}

Here one near-optimal point can use a flat slice in one bag and another
point can use a flat slice in another bag. Constant shares must reserve
tolerance for every bag at once. In this finite-state class, a lower size
estimate based on $\Psi$ cannot have a polynomial loss in $K$ whose degree
is independent of $d$: fixing $d$ larger than twice that degree makes
the ratio above diverge. This does not establish a counterexample for
the continuous growth model, nor a converse to
\cref{thm:allocation-pointwise}.

\subsection{Why the Euclidean radius matters}
\begin{proposition}[Dimensional loss for a flat bag]
\label{prop:allocation-flat}
Let $d\ge1$, $\alpha>0$, and $\eps>0$. Let there be one bag with
$K_r=V_r=\{1,\ldots,d\}$, $X=[0,1]^d$, $F=0$, and
$m_B=-\alpha q_{B,\{1,\ldots,d\}}$. Write
$v_d=\vol\{z\in\R^d:\norm z\le1\}$.
Every certificate with $\LB\ge-\eps$ has at least
$(\alpha/\eps)^{d/2}/v_d$ leaves. If $0<\eps\le\alpha/4$, then
\[
 \sup_{\eta\ge0}\Phi_\infty(\eps,\eta)
 \le2^{-d}\ceil{\sqrt{\alpha/(4\eps)}}^d
 \le2^{-d}(\alpha/\eps)^{d/2}.
\]
Thus the ratio of the necessary leaf count to this sup-norm allocation
bound is at least $2^d/v_d$.
\end{proposition}
\begin{proof}
The root inequality on a leaf $B$ implies $\alpha q_B(z)\le\eps$ for
every $z\in B$. Set $R=\sqrt{\eps/\alpha}$ and divide $B$ into its
$2^d$ nearest-vertex regions, cutting each coordinate at its midpoint.
The nearest-endpoint argument in
\cref{thm:allocation-pointwise} puts each region in one orthant of a
Euclidean ball of radius $R$ centered at its vertex. Each has volume at
most $v_dR^d/2^d$, so $\vol B\le v_dR^d$. The leaves cover a cube of
volume one, which proves the leaf bound.

For every $\eta$, $E(\eta)=X$, and the constant function $e=\eps$ is
admissible. Partition each coordinate interval into
$n=\ceil{1/(2R)}$ intervals of length at most $2R$. Their product centers
cover $X$ by $n^d$ sup-norm balls of radius $R$. This proves the first
upper bound. Since $1/(2R)\ge1$, the inequality
$\ceil{u}\le2u$ gives the second.
\end{proof}

The distinction is within a single bag, without separator inconsistency:
the lower gap adds squared coordinate distances, whereas a sup-norm
radius bounds only their maximum. The pointwise Euclidean bound retains
that sum. These necessary bounds give no general two-sided
characterization of the smallest affine separator certificate.
